\chapter[1927 --- Wagner-Jauregg]{1927: Julius Wagner-Jauregg}
\label{ch:1927_juliuswagnerjauregg}

\section*{Main Contribution}

\textit{Discovered that inoculation with malaria could cure general paralysis (tertiary syphilis), pioneering the use of fever therapy for infectious diseases.}

\section{Official Citation}

for his discovery of the therapeutic value of inoculation against general paralysis

\section{Background and Historical Context}

Julius Wagner-Jauregg was born on June 7, 1857, in Wels, Upper Austria. He studied medicine at the University of Vienna, where he was influenced by the distinguished psychiatrist Theodor Meynert and the neurologist Richard von Krafft-Ebing, both of whom were leading figures in the emerging field of neuropsychiatry. After receiving his medical degree in 1880, Wagner-Jauregg worked at the Vienna General Hospital, where he developed a strong interest in the treatment of mental illness and neurological disease. His early career was marked by a relentless search for effective therapies for conditions that were otherwise untreatable, and this search would eventually lead him to the development of pyrotherapy---the deliberate induction of fever as a therapeutic intervention.

Wagner-Jauregg's interest in fever as a therapeutic agent was inspired by a clinical observation that had been made by several physicians since the mid-nineteenth century: patients who contracted a febrile illness (such as erysipelas, pneumonia, or malaria) while suffering from certain chronic diseases sometimes experienced a temporary or permanent improvement in their symptoms. The most striking examples involved patients with general paresis of the insane---a progressive dementia caused by neurosyphilis---who seemed to improve when they developed a high fever from an intercurrent infection. Wagner-Jauregg hypothesized that the fever itself, rather than the specific infectious agent, was responsible for the improvement, and he set about testing this hypothesis through a series of experiments that would span more than two decades.

Wagner-Jauregg's Nobel Prize was awarded for a therapeutic intervention that was as controversial as it was significant: the deliberate inoculation of malaria parasites to treat neurosyphilis. To appreciate the significance of this achievement, it is necessary to understand the devastating impact of neurosyphilis on the populations of Europe and America in the late nineteenth and early twentieth centuries.

Neurosyphilis is a late complication of syphilis, a sexually transmitted infection caused by the spirochete bacterium \emph{Treponema pallidum}. After the initial infection, which is transmitted through sexual contact, syphilis progresses through several stages: a primary stage (characterized by a painless ulcer called a chancre at the site of infection), a secondary stage (characterized by a systemic rash, fever, and generalized lymphadenopathy), a latent stage (which can last for years or decades without symptoms), and a tertiary stage (in which the infection damages the heart, brain, and other organs). Neurosyphilis, the tertiary stage involving the central nervous system, manifests in several forms, the most severe of which is general paresis of the insane (also known as ``general paresis'' or ``paretic neurosyphilis'')---a progressive dementia characterized by personality changes, delusions of grandeur, memory loss, seizures, and ultimately death.

In the early twentieth century, general paresis was one of the most common causes of admission to psychiatric hospitals. It accounted for approximately 10--15\% of all admissions to mental institutions in Europe and North America, and the prognosis was uniformly fatal. There was no effective treatment, and patients typically died within three to five years of diagnosis. The wards of psychiatric hospitals were filled with patients suffering from this terrible disease---once proud and accomplished individuals reduced to confused, incontinent, and demented shells of their former selves---and the medical profession was desperate for any intervention that might offer hope.

\section{The Discovery and Contribution}

Wagner-Jauregg's path to the development of pyrotherapy was long and circuitous. He first attempted to induce fever in paretic patients by inoculating them with streptococci (the cause of erysipelas) in 1887. The results were mixed: some patients showed improvement, but others did not, and the streptococcal inoculations were difficult to control and potentially dangerous. Wagner-Jauregg then tried tuberculin (the tuberculin test material developed by Robert Koch for the diagnosis of tuberculosis), but this agent produced only modest and unreliable fevers.

It was not until 1917, when Wagner-Jauregg gained access to malaria-infected blood from soldiers returning from the Balkan front of World War I, that he found the agent he had been seeking. Malaria was an ideal pyrotherapeutic agent for several reasons: it produced high, sustained fevers (temperatures of 40--41\,$^{\circ}$C) that could be maintained for several days or weeks; the infection could be controlled with quinine, an effective antimalarial drug; and the fever pattern could be regulated by timing the administration of quinine.

On June 14, 1917, Wagner-Jauregg inoculated nine patients with general paresis with blood containing \emph{Plasmodium vivax}, the parasite that causes benign tertian malaria. The results were remarkable: within days of inoculation, the patients developed high fevers, and after several malarial paroxysms, many of them showed dramatic improvement in their mental symptoms. Some patients recovered completely, regaining their memory, their clarity of thought, and their ability to function in daily life. Others showed partial improvement that allowed them to leave the hospital and resume useful lives. A few patients did not respond, but even among non-responders, the progression of the disease was often slowed.

Wagner-Jauregg refined his technique over the following years, developing protocols for the optimal timing and dosage of malarial inoculation and for the subsequent treatment of the induced malaria with quinine. He showed that the therapeutic effect depended on the intensity and duration of the fever: patients who developed high fevers for a sufficient period were more likely to improve than those who experienced only mild or brief febrile episodes. He also demonstrated that the therapeutic effect was specific to malaria: other febrile infections were less effective, and artificial fever (induced by physical means such as hot baths or diathermy) was generally unsuccessful.

The mechanism by which malaria therapy improved neurosyphilis was not understood at the time, and it remains a matter of debate today. Several hypotheses have been proposed. The most widely accepted explanation is that the high fever induced by malaria enhanced the immune response against \emph{T. pallidum}, which thrives in the relatively cool environment of the central nervous system. The elevated body temperature may have directly killed the spirochetes or rendered them more susceptible to immune attack. An alternative hypothesis is that the malaria-induced immune activation produced a generalized immunostimulation that enhanced the clearance of the spirochetes from the brain. A third possibility is that the fever-induced increase in blood-brain barrier permeability allowed immune cells and antibodies to access the central nervous site of infection more effectively. Modern research has shown that fever does indeed enhance immune function by increasing the proliferation and activation of T lymphocytes, natural killer cells, and macrophages, providing a plausible mechanistic explanation for Wagner-Jauregg's clinical observations.

Whatever the mechanism, the clinical results were impressive. In Wagner-Jauregg's series of over 200 patients, approximately one-third showed complete recovery, another third showed significant improvement, and only one-third failed to respond. These outcomes were far superior to the natural history of untreated general paresis, in which virtually all patients died within a few years. The introduction of malaria therapy represented the first effective treatment for neurosyphilis, and it remained the treatment of choice until the development of penicillin in the 1940s.

Wagner-Jauregg received the Nobel Prize in Physiology or Medicine in 1927 ``for his discovery of the therapeutic value of inoculation of malaria paralysis in progressive paralysis.'' The award was not without controversy: some critics argued that deliberately infecting patients with a potentially fatal disease was unethical, and others questioned whether the benefits of malaria therapy truly outweighed its risks. Wagner-Jauregg responded to these criticisms by pointing out that the alternative---the untreated progression of general paresis to death---was far worse than the risks of induced malaria, which could be controlled with quinine.

\section{Impact on Modern Medicine}

Wagner-Jauregg's pyrotherapy represented a paradigm shift in the treatment of infectious disease: the idea that inducing a systemic physiological response (fever) could be therapeutic, even if the specific pathogen causing the fever was not directly targeting the disease-causing organism. This concept anticipated the modern understanding of fever as an immunomodulatory response and the use of immunotherapy in the treatment of infectious and neoplastic diseases.

The principle that fever enhances immune function is now well established. Studies have shown that elevated body temperature enhances the proliferation and activation of T lymphocytes, natural killer cells, and macrophages, and that the production of pro-inflammatory cytokines (such as interleukin-1, interleukin-6, and tumor necrosis factor) is enhanced at febrile temperatures. These findings provide a mechanistic explanation for the therapeutic effects that Wagner-Jauregg observed, and they have informed the development of modern immunotherapeutic strategies.

The concept of pyrotherapy has been revived in recent years in the context of cancer immunotherapy. Whole-body hyperthermia (the deliberate elevation of body temperature to 40--42\,$^{\circ}$C) is being investigated as an adjunct to chemotherapy and radiation therapy for various cancers, based on the principle that elevated temperature enhances immune recognition and destruction of tumor cells. Similarly, the use of Toll-like receptor agonists and other immunostimulatory agents to induce systemic inflammatory responses in cancer patients is conceptually related to Wagner-Jauregg's use of malaria-induced fever. The approval of immune checkpoint inhibitors (such as pembrolizumab and nivolumab) for the treatment of cancer has created a renewed interest in the mechanisms of immune activation that Wagner-Jauregg first exploited therapeutically.

More broadly, Wagner-Jauregg's work exemplifies the principle that medical progress sometimes requires bold and unconventional thinking. The idea of deliberately infecting patients with malaria to treat another disease was counterintuitive and controversial, but Wagner-Jauregg's clinical results were so compelling that the treatment was quickly adopted worldwide. His willingness to challenge conventional wisdom and to pursue a therapeutic hypothesis based on clinical observation, despite significant ethical and scientific objections, is a model for medical innovators.

The treatment of neurosyphilis was ultimately revolutionized by the introduction of penicillin in the 1940s, which provided a safe and effective cure for syphilis at all stages. Malaria therapy was gradually abandoned as penicillin became widely available, and it is no longer used today. However, the legacy of Wagner-Jauregg's work extends far beyond the specific treatment of neurosyphilis. His demonstration that the body's own physiological responses can be harnessed for therapeutic purposes was a precursor to the immunotherapy revolution that is transforming the treatment of cancer, autoimmune diseases, and infectious diseases in the twenty-first century.

\section{Legacy}

Julius Wagner-Jauregg's legacy is one of clinical courage and intellectual boldness. He was a physician who refused to accept the hopelessness of a fatal disease, and who had the imagination to envision a therapeutic approach that was radically different from anything that had been tried before. His work saved thousands of lives and demonstrated that even the most devastating diseases could be treated, if physicians were willing to think creatively and to accept calculated risks.

Wagner-Jauregg was also a thoughtful advocate for the ethical conduct of medical research. He recognized that his treatment carried risks, and he took pains to ensure that patients were fully informed and that the risks were minimized through careful monitoring and the timely use of quinine. His approach to the ethics of experimental therapy---balancing the potential benefits against the potential harms, and insisting on informed consent---anticipates the principles of modern research ethics that were formalized in the Declaration of Helsinki and the Belmont Report.

The Wagner-Jauregg Institute at the University of Vienna, which he founded, continues to conduct research in psychiatry and neurology. His name is commemorated in the Wagner-Jauregg Medal, awarded by the Austrian Society of Psychiatry and Neurology for outstanding contributions to the field. But perhaps the most enduring aspect of his legacy is the demonstration that fever---that ancient and universal response to infection---is not merely a symptom to be suppressed but a powerful physiological response that can be harnessed for therapeutic benefit. Wagner-Jauregg died in Vienna on May 27, 1940, but his vision of immunotherapy continues to inspire investigators who seek to mobilize the body's own defenses against disease.


\section{Modern Developments}

Wagner-Jauregg's fever therapy has been superseded by antibiotics but inspired modern immunotherapy. Hyperthermia therapy, using controlled heating to enhance chemotherapy effectiveness, is approved for certain cancers. Understanding of the immune-stimulating effects of fever has informed cancer immunotherapy approaches, including cytokine therapy (IL-2, interferon-alpha) and cancer vaccines. The concept of using infections to treat cancer has been revived with oncolytic virus therapy (T-VEC, approved 2015).
